\documentclass[12pt]{article}
\usepackage[utf8]{inputenc}
\usepackage[spanish]{babel}
\usepackage[a4paper, margin=2cm]{geometry}
\usepackage{tikz}
\usetikzlibrary{babel}
\usetikzlibrary{calc}
\usepackage[most]{tcolorbox}
\usepackage{amsmath}

% Usar fuente sin serifa para un aspecto más moderno (tipo infografía)
\renewcommand{\familydefault}{\sfdefault}

% Definición de colores con temática médica/científica
\definecolor{medblue}{RGB}{43, 108, 176}
\definecolor{medteal}{RGB}{49, 151, 149}
\definecolor{medlight}{RGB}{235, 248, 255}
\definecolor{meddark}{RGB}{26, 32, 44}
\definecolor{lungpink}{RGB}{244, 114, 182}  % Rosa claro para pulmones
\definecolor{lungdark}{RGB}{159, 18, 57}    % Rojo oscuro para bronquios/bordes

\begin{document}
\pagestyle{empty}

\begin{tcolorbox}[
    enhanced, 
    colback=medlight!30, 
    colframe=medblue, 
    title={\Large \textbf{Potencias de Diez: Cantidades Grandes}}, 
    halign title=center, 
    fonttitle=\bfseries, 
    arc=4mm, 
    boxrule=1.5pt,
    shadow={2mm}{-2mm}{0mm}{black!30!white}
]

    \vspace{0.2cm}
    \begin{tcolorbox}[colback=white, colframe=medteal, arc=2mm, boxrule=1.2pt, title=\textbf{Caso Clínico / Problema}]
        Los pulmones humanos contienen un promedio de \textbf{480,000,000} de alvéolos para el intercambio gaseoso. Escriba esta cifra empleando potencias de diez (notación científica).
    \end{tcolorbox}

    \vspace{0.4cm}
    
    \begin{center}
    \begin{tikzpicture}[scale=1]
        
        % Tráquea y bronquios
        \draw[line width=4pt, lungdark, cap=round] (0, 1.8) -- (0, 0.4);
        \draw[line width=3pt, lungdark, cap=round] (0, 0.4) -- (-0.8, -0.2);
        \draw[line width=3pt, lungdark, cap=round] (0, 0.4) -- (0.8, -0.2);
        
        % Pulmón izquierdo (dibujado esquemáticamente con elipse rotada)
        \filldraw[lungpink!40, draw=lungdark, line width=1.5pt, opacity=0.9, rotate around={15:(-1.2,-0.8)}] (-1.2,-0.8) ellipse (1cm and 1.8cm);
        % Pulmón derecho
        \filldraw[lungpink!40, draw=lungdark, line width=1.5pt, opacity=0.9, rotate around={-15:(1.2,-0.8)}] (1.2,-0.8) ellipse (1cm and 1.8cm);
        
        % Alvéolos ilustrativos (racimos dentro de los pulmones)
        \foreach \x/\y in {-1.2/-0.8, -0.8/-1.5, -1.6/-1.8, 1.2/-0.8, 0.8/-1.5, 1.6/-1.8} {
            \begin{scope}[shift={(\x,\y)}]
                \filldraw[white, draw=lungdark, thick] (0,0) circle (0.15cm);
                \filldraw[lungpink, draw=lungdark, very thin] (0,0) circle (0.06cm);
                \filldraw[lungpink, draw=lungdark, very thin] (0.1,0.06) circle (0.05cm);
                \filldraw[lungpink, draw=lungdark, very thin] (-0.05,0.1) circle (0.06cm);
                \filldraw[lungpink, draw=lungdark, very thin] (0.05,-0.08) circle (0.05cm);
            \end{scope}
        }
        
        \node[font=\bfseries, text=lungdark, fill=white, inner sep=2pt, rounded corners=2pt] at (0, 2.3) {Sistema Respiratorio};
        
        % Representación visual de la notación científica
        \node[font=\Large, text=meddark] at (0, -3.8) {
            $\text{N}^{\circ} = 4 \textcolor{lungdark}{\mathbf{,}} \underbrace{80\,000\,000}_{\text{\small \textbf{8 lugares a la izquierda}}} \text{ alvéolos} \quad \xrightarrow{\text{Notación}} \quad \mathbf{4.8 \times 10^8}$
        };
        
    \end{tikzpicture}
    \end{center}

    \vspace{0.4cm}
    \begin{tcolorbox}[colback=medteal!10, colframe=medteal, arc=2mm, boxrule=1.2pt, title=\textbf{Desarrollo Matemático y Solución}]
        Para convertir un número entero tan grande a notación científica, la coma decimal implícita al final del número debe desplazarse hacia la izquierda hasta formar un coeficiente comprendido entre el 1 y el 10.
        
        Al mover la coma \textbf{8 posiciones hacia la izquierda} (ubicándola entre el 4 y el 8), multiplicamos por una base 10 con un exponente positivo de \textbf{8}:
        \[ \text{N}^{\circ} \text{ Alvéolos} = \mathbf{4.8 \times 10^8} \]
        \textbf{Resultado:} Los pulmones humanos contienen en promedio \textbf{$4.8 \times 10^8$} alvéolos.
    \end{tcolorbox}
    
\end{tcolorbox}

\end{document}